\section*{Capítulo \ref{ch:b2:reaction-enthalpy} --- Entalpías de reacción}

\begin{solution}{exo:b2:reaction-enthalpy:1}
$\Delta_r H^\circ = 3(-393.51) + 4(-285.83) - (-104.7) - 5(0) =
\qty{-2219.2}{kJ/mol}$, negativa: \omterm{def:b2:reaction-enthalpy:exothermic}{exotérmica}. (La entalpía de combustión
medida es la misma, \qty{-2219.2}{kJ/mol}.)
\end{solution}

\begin{solution}{exo:b2:reaction-enthalpy:2}
Metano: $890.3/16.04 = \qty{55.5}{kJ/g}$, es decir, \qty{55.5}{MJ} por
kilogramo. Dihidrógeno: $285.83/2.016 = \qty{141.8}{kJ/g}$,
\qty{141.8}{MJ/kg}: unas 2.6 veces más por kilogramo (pero mucho menos por
litro, por ser un gas ligero).
\end{solution}

\begin{solution}{exo:b2:reaction-enthalpy:3}
$\Delta\nu_{\text{gas}} = 6 - 7.5 = -1.5$, luego $\Delta_r H = \Delta_r U +
\Delta\nu_{\text{gas}}RT = -3264 - 1.5 \times 8.314 \times 298.15 \times
10^{-3} = \qty{-3267.7}{kJ/mol}$, de acuerdo con el valor calculado a partir de
las entalpías de formación, \qty{-3267.6}{kJ/mol}.
\end{solution}

\begin{solution}{exo:b2:reaction-enthalpy:4}
La reacción buscada es la primera menos dos veces la segunda:
$\Delta_r H^\circ = -393.5 - 2(-283.0) = \qty{+172.5}{kJ/mol}$. Es
\omterm{def:b2:reaction-enthalpy:exothermic}{endotérmica}: el coque caliente se enfría al convertir el dióxido de carbono en monóxido
de carbono.
\end{solution}

\begin{solution}{exo:b2:reaction-enthalpy:5}
Enlaces rotos: C--H y Cl--Cl, $416 + 243 = \qty{659}{kJ/mol}$; enlaces
formados: C--Cl y H--Cl, $350 + 432 = \qty{782}{kJ/mol}$: estimación
\qty{-123}{kJ/mol}. A partir de las entalpías de formación: $-83.7 - 92.3 + 74.9 =
\qty{-101.1}{kJ/mol}$. La diferencia de \qty{22}{kJ/mol} procede del valor medio de C--Cl,
promedio sobre varias moléculas, y del enlace C--H roto en el
metano, más fuerte que la media de los cuatro.
\end{solution}

\begin{solution}{exo:b2:reaction-enthalpy:6}
A \qty{298}{K}: $-241.83 + 285.83 = \qty{44.00}{kJ/mol}$. Kirchhoff con
$\Delta C_p = 33.59 - 75.35 = \qty{-41.76}{J/(K.mol)}$:
$44.00 - 0.04176 \times 101.85 = \qty{39.75}{kJ/mol}$ a \qty{400}{K}. Las
tablas dan $-242.85 + 282.59 = \qty{39.74}{kJ/mol}$: las capacidades
caloríficas constantes son excelentes en este intervalo.
\end{solution}

\begin{solution}{exo:b2:reaction-enthalpy:7}
Ionización: $4.341 \times 96.485 = \qty{418.8}{kJ/mol}$; captura:
$-3.613 \times 96.485 = \qty{-348.6}{kJ/mol}$. \omterm{def:b2:reaction-enthalpy:lattice-enthalpy}{Entalpía reticular} $= 436.7 +
89.0 + 121.3 + 418.8 - 348.6 = \qty{717}{kJ/mol}$, menor que los
\qty{787}{kJ/mol} del cloruro de sodio: \ce{K+} es mayor que \ce{Na+}, los
iones están más alejados y se atraen menos.
\end{solution}

\begin{solution}{exo:b2:reaction-enthalpy:8}
Capacidad calorífica del calorímetro: $C = 2.00/0.418 = \qty{4.785}{kJ/K}$.
Calor liberado: $4.785 \times 0.583 = \qty{2.789}{kJ}$, para $\xi =
\qty{0.0500}{mol}$ (el ácido es el reactivo limitante, la base está en ligero exceso):
$\Delta_r H = -2.789/0.0500 = \qty{-55.8}{kJ/mol}$.
\end{solution}

\begin{solution}{exo:b2:reaction-enthalpy:9}
\ce{C8H18(l) + 25/2O2(g) -> 8CO2(g) + 9H2O(l)}: $\Delta_r H^\circ =
8(-393.51) + 9(-285.83) + 250.3 = \qty{-5470.3}{kJ/mol}$; con $M =
\qty{114.2}{g/mol}$, \qty{47.9}{MJ/kg}. Dióxido de carbono: $8 \times 44.01 =
\qty{352.1}{g}$ por \qty{5.470}{MJ}, es decir, \qty{64.4}{g/MJ}, frente a
$44.01/0.8903 = \qty{49.4}{g/MJ}$ para el metano: el metano tiene más hidrógeno por
carbono, y la combustión del hidrógeno libera calor sin dióxido de carbono.
\end{solution}

\begin{solution}{exo:b2:reaction-enthalpy:10}
Productos de \ce{CH4 + 2O2 -> CO2 + 2H2O(g)}: $C_p = 37.13 + 2 \times 33.59 =
\qty{104.3}{J/K}$; $q = \qty{802.3}{kJ}$; $T_f = 298 + 802\,300/104.3 \approx
\qty{7990}{K}$, casi el triple del valor en aire, porque no hay nitrógeno
que se lleve su parte del calor. La llama real en oxígeno es mucho más fría: mucho antes de \qty{7990}{K}, el dióxido de carbono y el agua
se disocian en monóxido de carbono, hidrógeno, oxígeno y átomos, reacciones
\omterm{def:b2:reaction-enthalpy:exothermic}{endotérmicas} que absorben la mayor parte del calor; además, las capacidades caloríficas crecen con
$T$.
\end{solution}

\begin{solution}{exo:b2:reaction-enthalpy:11}
Se obtiene $\Delta_r H^\circ(700) = \Delta_r H^\circ(298) + \int_{298}^{700}(\Delta_r a +
\Delta_r b\,T)\,\dd T = -91.88 + [-60.5 \times 401.85 + 0.0270 \times
(700^2 - 298.15^2)] \times 10^{-3} = -91.88 - 24.31 + 10.83 =
\qty{-105.4}{kJ/mol}$, a \qty{0.2}{kJ/mol} de las tablas
(\qty{-105.2}{kJ/mol}), cuando la estimación con $C_p$ constante se desviaba
\qty{4.5}{kJ/mol}.
\end{solution}

\begin{solution}{exo:b2:reaction-enthalpy:12}
$\Delta_r H^\circ = 2 \times 90.29 = \qty{+180.6}{kJ/mol}$ (\omterm{def:b2:reaction-enthalpy:exothermic}{endotérmica}).
$\Delta_r C_p^\circ = 2(29.85) - 29.12 - 29.38 = \qty{1.2}{J/(K.mol)}$, luego
entre \qty{298}{K} y \qty{2000}{K} la \omterm{def:b2:reaction-enthalpy:reaction-quantity}{entalpía de reacción} varía en
unos $1.2 \times 1702 = \qty{2}{kJ/mol}$, un uno por ciento. Una reacción \omterm{def:b2:reaction-enthalpy:exothermic}{endotérmica}
se ve favorecida por la temperatura alta (\cref{ch:b2:equilibrium-shifts}):
el aire frío apenas produce monóxido de nitrógeno, y las llamas más calientes son las que
más producen; por eso la temperatura de combustión es una palanca contra este
contaminante.
\end{solution}

\begin{solution}{pb:b2:reaction-enthalpy:1}
\textbf{1.} \ce{C4H10(g) + 13/2O2(g) -> 4CO2(g) + 5H2O(l)}.
\textbf{2.} $\Delta_r H^\circ = 4(-393.51) + 5(-285.83) + 125.6 =
\qty{-2877.6}{kJ/mol}$.
\textbf{3.} Con vapor de agua, $4(-393.51) + 5(-241.83) + 125.6 =
\qty{-2657.6}{kJ/mol}$; la diferencia, \qty{220.0}{kJ}, es la entalpía de
vaporización de cinco moles de agua (\qty{44.0}{kJ/mol} cada uno).
\textbf{4.} $2877.6/58.12 = \qty{49.5}{kJ/g}$ y $2657.6/58.12 =
\qty{45.7}{kJ/g}$.
\textbf{5.} El segundo: el agua sale como vapor y su calor de
condensación se pierde.
\textbf{6.} $230 \times 45.7 = \qty{1.05e4}{kJ}$, unos \qty{10.5}{MJ}.
\textbf{7.} A volumen constante el calor recibido es $\Delta U$ (no hay
trabajo de las fuerzas de presión): la bomba mide $\Delta_r U$.
\textbf{8.} $\Delta\nu_{\text{gas}} = 4 - 1 - 6.5 = -3.5$ (el agua es
líquida).
\textbf{9.} $\Delta_r U = \Delta_r H - \Delta\nu_{\text{gas}}RT = -2877.6 +
3.5 \times 2.479 = \qty{-2868.9}{kJ/mol}$.
\textbf{10.} $\xi = 0.500/58.12 = \qty{8.60e-3}{mol}$; calor $8.60 \times
10^{-3} \times 2868.9 = \qty{24.68}{kJ}$; $\Delta T = 24.68/10.00 =
\qty{2.468}{K}$.
\textbf{11.} Calor $10.00 \times 2.461 = \qty{24.61}{kJ}$: $\Delta_r U =
-24.61/(8.603 \times 10^{-3}) = \qty{-2860.6}{kJ/mol}$ y $\Delta_r H =
-2860.6 - 8.7 = \qty{-2869.3}{kJ/mol}$, un \qty{0.29}{\percent} por debajo de las
tablas en valor absoluto: calor perdido por el calorímetro o una combustión
ligeramente incompleta.
\textbf{12.} $n = 1000/18.02 = \qty{55.5}{mol}$; $Q = 55.5 \times 75.35
\times 80 = \qty{334.5}{kJ}$.
\textbf{13.} Calor liberado $16.0 \times 45.72 = \qty{731.5}{kJ}$; rendimiento
$334.5/731.5 = 0.46$.
\textbf{14.} A los gases de combustión calientes que escapan alrededor del cazo, al
aire calentado por la llama y a las paredes del cazo, y a la radiación.
\textbf{15.} $230 \times 45.72 \times 0.457/334.5 = \qty{14.4}{L}$.
\textbf{16.} Los \qty{6.5}{mol} de \ce{O2} vienen con $6.5 \times 78.08/20.95 =
\qty{24.23}{mol}$ de \ce{N2} y $6.5 \times 0.93/20.95 = \qty{0.289}{mol}$
de \ce{Ar}.
\textbf{17.} \qty{4}{mol} de \ce{CO2}, \qty{5}{mol} de \ce{H2O(g)},
\qty{24.23}{mol} de \ce{N2} y \qty{0.289}{mol} de \ce{Ar}.
\textbf{18.} La llama es adiabática e isobara, luego $\Delta H = 0$. A lo largo de un
camino formado por la reacción a \qty{298}{K} ($\Delta H_1 =
\qty{-2657.6}{kJ}$) seguida del calentamiento de estos productos, $\Delta H_2 =
-\Delta H_1 = \qty{+2657.6}{kJ}$: todo el calor va a los gases.
\textbf{19.} $C_p = 4(37.13) + 5(33.59) + 24.23(29.12) + 0.289(20.79) =
\qty{1028}{J/K}$.
\textbf{20.} $T_f = 298 + 2\,657\,600/1028 = \qty{2883}{K}$.
\textbf{21.} $T_f = 2300 + 100 \times (2657.6 - 2515.7)/(2655.7 - 2515.7) =
\qty{2401}{K}$.
\textbf{22.} Las capacidades caloríficas de los gases crecen con la temperatura (en
torno a una cuarta parte entre \qty{298}{K} y \qty{2400}{K}): los valores a temperatura
ambiente calientan los gases demasiado barato.
\textbf{23.} Por encima de \qty{2000}{K} el dióxido de carbono y el agua se disocian
en parte (de forma \omterm{def:b2:reaction-enthalpy:exothermic}{endotérmica}), la llama radia y la mezcla con el aire nunca es perfecta.
\textbf{24.} La \omterm{def:b2:reaction-enthalpy:flame-temperature}{temperatura adiabática de llama} del butano en aire es
$\boldsymbol{T_{\text{ad}} \approx \qty{2.40e3}{K}}$.
\end{solution}
